\begin{textsample}{POS Dim 3 – to – Score 24.00 – to/to\_inf\_29.txt}  \label{ex:f3_pos_106}
O Campo de Experiência “ traços, \textbf{sons}, cores e formas ” O Campo de Experiência “ traços, \textbf{sons}, cores e formas ” propõe desenvolver e valorizar as diferentes linguagens e manifestações artísticas, culturais, simbólicas e científicas, relacionadas aos contextos sociais em que as \textbf{crianças} estão inseridas.
Considerando a \textbf{criança} como um ser histórico e social, deve -se olhar para o \textbf{lúdico} como um precioso recurso, deixando assim, a aprendizagem mais significativa, pois o \textbf{brincar} e o interagir fazem parte do mundo \textbf{infantil}.
Campo de experiência : traços, \textbf{sons}, cores e formas As \textbf{crianças} constituem sua identidade pessoal e social nas interações com os demais atores sociais, durante as quais elas se apropriam, descobrem e se expressam por meio das variadas linguagens, no contato com manifestações culturais locais e de outros países.
Por isso, é importante que, desde \textbf{bebês}, as \textbf{crianças} tenham momentos para \textbf{manusear} diferentes materiais, realizando suas produções com \textbf{gestos}, \textbf{sons}, danças, \textbf{mímicas}, traços, encenações, desenhos, modelagens, canções, de modo singular, inventivo e prazeroso, desenvolvendo sua sensibilidade.
Quando o professor abordar os artistas com as \textbf{crianças}, não focar primordialmente em dados biográficos, o objetivo não é decorarem nomes ou copiarem obras.
E sim, experimentar os processos expressivos, conhecer os materiais, explorar, imaginar e \textbf{inventar}.
As práticas pedagógicas precisam permitir à \textbf{criança} a liberdade de escolha, respeitar o tempo de cada uma, levar em conta o que elas têm a dizer, quer seja com o corpo ou com a voz.
Ações que o professor conduz para que a \textbf{infância} seja verdadeiramente vivida, explorada e experimentada em todas as suas possibilidades.
A observação \textbf{atenta} do \textbf{adulto} com relação à \textbf{criança} vem ao encontro de uma postura que a valoriza em sua \textbf{curiosidade} e \textbf{desejo} de conhecer o mundo.
Olhar a \textbf{criança} para além do que os olhos veem, prestar atenção aos detalhes do que produzem, pensam, expressam, desejam e interagem.
É importante quando as \textbf{crianças} transformam o material naquilo que desejam, quando blocos de \textbf{montar} viram \textbf{personagens}, num diálogo \textbf{cheio} de significâncias.
Para além de grafar somente a pontinha do \textbf{dedo} na tinta, acredita -se em propostas em que a \textbf{criança} possa conhecer o mundo, por meio do próprio corpo e das múltiplas linguagens, descobrindo modos próprios de se expressar, ampliando assim seu repertório.
As \textbf{crianças} elaboram a sua própria organização e criação, e necessitam se envolver com propostas atrativas para aprender e construir conhecimentos.
Porque é na interação com os \textbf{pares}, e principalmente entre as próprias \textbf{crianças}, que são as verdadeiras artistas propositoras, que o processo de desenvolvimento será fortalecido.
É preciso ofertar às \textbf{crianças} caminhos com experiências significativas para \textbf{infâncias} criativas, brincantes e livres.
Os campos de experiências norteiam as atividades, porém, toda ação deve estar centralizada na \textbf{criança} e seus interesses como pesquisadora, favorecendo as manifestações culturais mais significativas, materiais e tecnologias, realizando produções com \textbf{gestos}, traços, desenhos, modelagens, danças, jogos simbólicos, \textbf{sons} e canções.
A Educação \textbf{Infantil} precisa promover a participação das \textbf{crianças} em atividades de produção, manifestação e apreciação da arte, de modo a favorecer o desenvolvimento da sensibilidade, ludicidade, criatividade e das diversas expressões artísticas e culturais, permitindo assim, que elas potencializem suas singularidades, as relações e vivências artísticas.
Traços, \textbf{sons}, formas e imagens - CONVIVER e fruir das manifestações artísticas e culturais da sua comunidade e de outras culturas - artes \textbf{plásticas}, música, dança, teatro, cinema, folguedos e festas populares - ampliando a sua sensibilidade, desenvolvendo senso estético, empatia e respeito às diferentes culturas e identidades. - \textbf{BRINCAR} com diferentes \textbf{sons}, ritmos, formas, cores, \textbf{texturas}, objetos, materiais, construindo cenários e indumentárias para \textbf{brincadeiras} de faz de conta, encenações ou para festas tradicionais, enriquecendo seu repertório e desenvolvendo seu senso estético. - PARTICIPAR de decisões e ações relativas à organização do ambiente ( tanto no cotidiano como na preparação de eventos especiais ), a definição de temas e a escolha de materiais a serem usados em atividades \textbf{lúdicas} e teatrais, entrando em contato com manifestações do patrimônio cultural, artístico e tecnológico, apropriando -se de diferentes linguagens. - EXPLORAR variadas possibilidades de usos e combinações de materiais, substâncias, objetos e recursos tecnológicos para criar e recriar danças, artes visuais, encenações teatrais, músicas, escritas e mapas, apropriando -se de diferentes manifestações artísticas e culturais. - EXPRESSAR, com criatividade e responsabilidade, suas emoções, sentimentos, necessidades e ideias \textbf{brincando}, cantando, dançando, esculpindo, desenhando, encenando, compreendendo e usufruindo o que é comunicado pelos demais \textbf{colegas} e pelos \textbf{adultos}. - CONHECER -SE, no contato criativo com manifestações artísticas e culturais locais e de outras comunidades, identificando e valorizando o seu pertencimento étnico racial, de gênero e de crença religiosa, desenvolvendo sua sensibilidade, criatividade, gosto pessoal e modo peculiar de expressão por meio do teatro, música, dança, desenho e imagens.
Brasil ( 2016 2ª Versão da BNCC )

% matched lemmas: adulto, atento, bebê, brincadeira, brincar, cheio, colega, criança, curiosidade, dedo, desejo, gesto, infantil, infância, inventar, lúdico, manusear, montar, mímico, par, personagem, plástico, som, textura
\end{textsample}
